\section{Methods}

\subsection{Autonomous-agent workflow, model and human oversight}

The autonomous programme ran under one continuous analysis protocol from 27 July to
14 August 2026. The multi-agent loop coordinated task assignment, parallel execution and
result exchange but did not itself call a large language model (LLM). Instead, our agents did all their reasoning by calling Codex directly through our own
\texttt{codex-api} client, which we built on the official OpenAI Python SDK
(\url{https://github.com/szl666/codex-api}). Codex used GPT-5.6-sol as its base model.
Across multiple sessions, our agents proposed hypotheses, designed and ran analyses, wrote
code, constructed checks and diagnosed failed claims in a shell with data access.

Humans set the initial goal, redirected the study scope and authorised access to data
reserved for later tests, but supplied no detailed scientific hypothesis. Confirmation
analyses followed written procedures fixed before evaluation. Human authors verified the
final outputs and retain responsibility for the reported conclusions. Agent records,
eleven refuted conclusions and reproduction commands are reported in Supplementary
Sections~\ref{si-note:s1}, \ref{si-note:s6}, \ref{si-note:s9}, \ref{si-note:s14}
and~\ref{si-note:s19}. Archived scripts and feature tables reproduce the reported values.

\subsection{First-principles verification}

Four quantities on which the analysis depends were initially learned rather than computed:
the contact thresholds of Law~1 and Law~2, the ordering behind GNoME's low-symmetry
excess, the damage-severity scale assigned by a machine-learned potential and the property
target used to select the design queue. Each was re-derived with plane-wave DFT under a protocol frozen
before any job was submitted. Calculations used VASP 6.3.0 with PBE PAW potentials, a
plane-wave cutoff of \(\max(520\,\mathrm{eV},1.3\max\mathrm{ENMAX})\), explicit
\(\Gamma\)-centred meshes at \(0.22\,\text{\AA}^{-1}\),
\(\mathrm{EDIFF}=10^{-6}\)\,eV and tetrahedron smearing.

The four tests measured the energy landscape along \rhoc{}, ordering energies over every
symmetry-distinct merge-group configuration, relaxation-energy release for matched cells,
and bulk moduli from third-order Birch--Murnaghan fits to five-point energy--volume curves.
Cells were fully relaxed before the corresponding static or constant-volume calculations,
and machine-learning and DFT relaxation energies used the same per-cell definition. The
campaign comprised 1{,}917 tasks. Supplementary Section~\ref{si-note:dft} gives the complete
protocol, per-quantity settings, convergence and fit-sensitivity analyses. Relaxed cells
for the 260 design candidates are provided as Supplementary Data.

\subsection{Data sets and study design}

Before evaluation, written protocols fixed the allowed structural quantities, data splits
and success criteria. We analysed 99{,}162 experimental structures from the Inorganic
Crystal Structure Database (ICSD) and the Crystallography Open Database (COD)
\cite{zagorac2019recent,grazulis2009crystallography}. For law discovery, a seeded hash of
each structure identifier assigned eligible ionic structures to discovery, held-out or
reserve partitions. Every damaged structure inherited its experimental parent's assignment.
The split therefore separated structures rather than chemistries: two structures with the
same reduced composition could fall in different partitions.

Thresholds were fitted on 12{,}632 experimental and 8{,}590 damaged discovery structures,
then assessed on 5{,}297 experimental and 3{,}612 damaged held-out structures without
refitting. We additionally evaluated the frozen law sets on held-out reduced compositions
absent from discovery. Damaged structures inherited their parent's composition, and
uncertainty was clustered by reduced composition. After law-set selection, a split-labelling error
allowed reserve structures into one full-sample threshold fit, so the reserve no longer
constitutes a fully independent final test. Database inventories, licences, the
composition-unseen analysis and the reserve incident are reported in Supplementary
Sections~\ref{si-note:s2}, \ref{si-note:s8} and~\ref{si-note:s19}.

\subsection{Structural descriptors and law-set evaluation}

We computed every Law~1--Law~8 quantity under a fixed convention. Formal oxidation states
were assigned from composition by integer charge balancing \cite{jablonka2021using}. Where
noted, external analyses used a non-integer mean-valence fallback without refitting any
law. Native CIF oxidation annotations and bond-length-based valence inference were excluded
to prevent bond lengths from entering both charge assignment and bond-valence evaluation.
Primary neighbours came from CrystalNN \cite{zimmermann2020local} in pymatgen
\cite{ong2013python}, and Law~7 used spglib \cite{togo2024spglib} with
\texttt{symprec}=0.01. Discovery used deposited cells, whereas external comparisons
involving Law~7 used a common primitive-cell convention.

Each law combines a structural quantity with a directional threshold and, where needed, an
explicit chemical condition. For benchmark rates, an unavailable measurement counted as
satisfying that law under the frozen convention, and only features available for more than
90\% of experimental structures were admitted to the search. In deployment, a missing
charge, radius or other required input instead produces a separate no-verdict outcome. Any
violation among the evaluable laws marks the structure implausible. Full descriptor
definitions, radius and bond-valence lookups, alternative-neighbour checks, missing-input
rules, cell conventions and the equations for Law~1--Law~8 are in Supplementary
Sections~\ref{si-note:s3}, \ref{si-note:s17} and~\ref{si-note:s18}.

\subsection{Chemically damaged structures and law selection}

The damaged structures came from five perturbations that preserved composition,
stoichiometry and atom count: uniaxial compression, cation--cation exchange, random
displacement, isotropic expansion and cation--anion exchange. Independent contact,
electrostatic and MatterSim-relaxation checks measured their physical and energetic
severity. We selected laws to maximise damage detection subject to an experimental-structure
satisfaction floor. Fixed Set~4 was then evaluated once on held-out data and compared with
a provably optimal decision tree within the tested depth-three class. A separate
threshold-transfer analysis re-derived each continuous cutoff at the corresponding held-out
percentile and measured the resulting verdict changes.

For leave-one-damage-class-out evaluation, tree and single-threshold selection were repeated
completely. For Set~4 we reselected only its additions on a Set~3 base that had already seen all five
classes, so this variant tests transfer of the added mechanisms, not de novo discovery. The distance-cutoff benchmark compared PRIS with fixed 0.5- and 0.7-\AA{}
cutoffs and one matched to Set~4 satisfaction, on 440 experimental structures
and 2{,}024 composition-preserving damaged variants. Exact perturbation operators,
search grids, physical checks, omitted-class boundaries and threshold-transfer results
are in Supplementary Sections~\ref{si-note:s4}, \ref{si-note:s5}, \ref{si-note:s7},
\ref{si-note:s10}, \ref{si-note:s11}, \ref{si-note:s12}, \ref{si-note:s17}
and~\ref{si-note:s18}.

\subsection{External evaluation, synthesizability screening and statistics}

With PRIS thresholds fixed, external evaluation covered GNoME and seven crystal generators,
similar-element merging, MatterSim relaxation, coordinate-based checks and historical
falsified depositions. Following a protocol fixed before computation, we sampled 5{,}000 of the 554{,}054
released GNoME structures. Separate analyses relaxed 500 raw MatterGen outputs and compared
thermodynamic stability, dynamical stability and experimental record across 26{,}600
Materials Project structures. Same-composition ranking covered 18{,}920 pairs spanning
1{,}508 compositions, weighted compositions equally, excluded ties from binary accuracy and
used composition-cluster bootstrap intervals.

PSS was fitted on development compositions as an antisymmetric, zero-intercept logistic
score of the six standardised descriptors in equation~\ref{eq:pss}. Unavailable descriptors
took frozen development-set medians. The score was evaluated once on held-out compositions
before transfer. To test synthesizability trends, we trained a 50-bag CGCNN-PU model on
99{,}162 experimental and 8{,}125{,}976 unlabelled structures. We independently trained 50 MLP PU
heads on frozen 128-dimensional MatterSim-v1.0.0-1M representations
\cite{jang2020structure,song2025llm,yang2024mattersim}. Consensus between the two models defined the
hard-to-synthesize cohort. For inverse design, 13 independently seeded MatterGen runs
conditioned at 400\,GPa yielded 1{,}081 unique candidates. Independent UMA bulk-modulus
predictions defined the high-property subset, and 541 experimental high-property structures
fixed the screening threshold, using the two PSS descriptors available for the candidates. Supplementary
Sections~\ref{si-note:s13}, \ref{si-note:s15}, \ref{si-note:pu-transfer},
\ref{si-note:s17} and~\ref{si-note:s19} give the data sets, descriptors, imputation,
coefficients, model validation, phase-hull construction, inverse-design checks and
statistical procedures.
